\documentclass[10pt]{article}
\usepackage[T1]{fontenc}
\usepackage[utf8]{inputenc}
\usepackage[french]{babel}
\usepackage{amsmath,amssymb,mathrsfs,mathpazo}
\usepackage{geometry}
\geometry{a4paper,margin=2cm}
\usepackage{graphicx,array,tabularx,tabularray,colortbl,multirow,multicol,booktabs,enumitem,pifont,ulem,fancybox,diagbox,cancel}
\usepackage{pstricks,pstricks-add,pst-plot,pst-tree,pst-node,pst-text,pst-func,pst-math,pst-bspline,pst-3dplot,pst-eucl,pst-solides3d}
\let\pscalculate\undefined
\usepackage{pst-bezier}
\usepackage{tikz,pgfplots,tkz-tab}
\usetikzlibrary{arrows,arrows.meta,patterns,calc,positioning,trees,decorations.markings,shapes.geometric,shapes.misc,shapes.arrows,intersections,decorations.pathmorphing,decorations.pathreplacing,backgrounds,shadows,babel}
\pgfplotsset{compat=1.18}
\usepgfplotslibrary{fillbetween}
\usepackage[np]{numprint}
\usepackage{listings,makecell,pas-tableur,twoopt}
\let\barre\undefined
\usepackage[pstricks]{tablvar}
\definecolor{darkgreen}{RGB}{0,100,0}
\providecommand{\red}{\color{red}}
\providecommand{\blue}{\color{blue}}
\providecommand{\green}{\color{green}}
\providecommand{\black}{\color{black}}
\providecommand{\white}{\color{white}}
\providecommand{\gray}{\color{gray}}
\providecommand{\cyan}{\color{cyan}}
\providecommand{\magenta}{\color{magenta}}
\providecommand{\yellow}{\color{yellow}}
\providecommand{\violet}{\color{violet}}
\providecommand{\orange}{\color{orange}}
\providecommand{\euro}{\texteuro}
\providecommand{\EUR}{\texteuro}
\providecommand{\ieme}{\textsuperscript{e}}
\providecommand{\degres}{\ensuremath{^\circ}}
\providecommand{\degree}{\ensuremath{^\circ}}
\providecommand{\vv}[1]{\overrightarrow{#1}}
\providecommand{\Coord}[1]{#1}
\providecommand{\decosix}{}
\providecommand{\starredbullet}{\ensuremath{\bullet}}
\pagestyle{empty}
\setlength{\parindent}{0pt}
\usepackage{adjustbox}\DeclareUnicodeCharacter{2236}{:}\DeclareUnicodeCharacter{2190}{\ensuremath{\leftarrow}}
\providecommand{\R}{\mathbb{R}}
\providecommand{\N}{\mathbb{N}}
\providecommand{\Z}{\mathbb{Z}}
\providecommand{\Q}{\mathbb{Q}}
\providecommand{\D}{\mathbb{D}}
\providecommand{\C}{\mathbb{C}}
\providecommand{\e}{{\mathrm{e}}}
\providecommand{\vect}[1]{\overrightarrow{#1}}
\providecommand{\vectt}[1]{\overrightarrow{\mathrm{#1}}}
\providecommand{\vvt}[1]{\overrightarrow{\mathrm{#1}}}
\providecommand{\np}[2][]{{#2~#1}}
\providecommand{\vv}[1]{{\overrightarrow{#1}}}
\renewcommand{\d}{{\mathrm{d}}}
\providecommand{\degres}{{^\circ}}
\providecommand{\degre}{{^\circ}}
\providecommand{\blue}{\color{blue}}
\providecommand{\red}{\color{red}}
\providecommand{\enskip}{\quad}
\providecommand{\llbracket}{\lbrack\!\lbrack}
\providecommand{\rrbracket}{\rbrack\!\rbrack}
\providecommand{\euro}{{\text{€}}}
\providecommand{\no}{{\text{n°}}}
\providecommand{\interval}[2]{{[#1;#2]}}
\providecommand{\intervFO}[2]{{[#1;#2[}}
\providecommand{\dotfill}{\ldots}

\providecommand{\ds}{\displaystyle}
\providecommand{\textding}[1]{#1}
\providecommand{\Oij}{\(\mathrm{O};\vec{\imath},\vec{\jmath}\)}
\providecommand{\Oijk}{\(\mathrm{O};\vec{\imath},\vec{\jmath},\vec{k}\)}
\providecommand{\Ouv}{\(\mathrm{O};\vec{u},\vec{v}\)}
\providecommand{\intervFF}[2]{\ensuremath{[#1;#2]}}
\renewcommand{\euro}{\ifmmode\text{\texteuro}\else\texteuro\fi}
\DeclareUnicodeCharacter{00B0}{\ifmmode^{\circ}\else\textdegree\fi}
\DeclareUnicodeCharacter{20AC}{\ifmmode\text{\texteuro}\else\texteuro\fi}
% Énoncés : Éducation nationale. Transcription et corrigés : APMEP, Denis Vergès et contributeurs. Découpage et classement : CoopMaths, Sébastien Lozano, Éric Elter et contributeurs. Adaptation de la présentation : Maths Explorer.
% CC BY-SA : https://coopmaths.fr/www/about/
% https://coopmaths.fr/alea/static/bac/2026/tex/bac_2026_09_sujet2_metropole_2_cor.tex
\begin{document}
\noindent\textit{APMEP / CoopMaths. CC BY-SA. Présentation adaptée par Maths Explorer.}\par\medskip
\begin{enumerate}
		\item On a :\quad
		$\displaystyle \lim_{x \to -\infty} 3x-1 = - \infty$
		\quad et \quad
		$\displaystyle \lim_{x \to -\infty} \e^{-x} = + \infty$,

		par limite du produit :\quad $\displaystyle \lim_{x \to -\infty} f(x) = - \infty$.

		\item On a :\quad
		$\displaystyle \lim_{x \to +\infty} 3x-1 = + \infty$
		\quad et \quad
		$\displaystyle \lim_{x \to +\infty} \e^{-x} = 0$, \quad sous cette forme, on a donc une forme indéterminée.

		Mais, pour tout $x$ réel, on a :\quad $f(x) = (3x- 1)\e^{-x} +2 =  -3 \big(-x \big)\e^{-x} - \e^{-x} +2$.


		$\displaystyle  \lim_{x \to +\infty} -x \e^{-x} = \displaystyle  \lim_{X \to -\infty} X \e^{X} = 0  $, \quad d'après la propriété des croissances comparées.

		On obtient alors, par limite de la somme :\quad $\displaystyle  \lim_{x \to +\infty} f(x) =2$.

		On en déduit que la droite d'équation $y=2$ est une asymptote horizontale à la courbe $\mathcal{C}_f$ en $+\infty$.

		\item \begin{enumerate}
			\item On a, pour tout réel $x$ : \quad $f'(x)=3 \times \e^{-x} + (3x-1) \times \left( - \e^{-x} \right)$.

			Soit \quad $f'(x)=\e^{-x} \big( 3-(3x-1) \big)=(-3x+4) \e^{-x}$.

			\item Pour tout $x \in \R$, \quad $\e^{-x} >0$ \quad donc le signe de $f'$ est le même que celui de $-3x+4$.

			Or \quad $-3x+4 \geqslant 0 \iff x \leqslant \dfrac43$.

			De plus, on a : \quad $f \left(\dfrac43 \right) = 3 \e^{-\frac43}+2$.

			On obtient le tableau de variations
			\begin{center}
				\begin{tikzpicture}
					\tkzTabInit[lgt=3,espcl=4]{$x$ / 0.8 , signe de $f'(x)$ / 1, variation de $f$ / 2}{$-\infty$, $\frac43$, $+\infty$}
					\tkzTabLine{,+,z,-, }
					\tkzTabVar{-/$-\infty$, +/$3\e^{\frac43}+2$,  -/$2$ }
				\end{tikzpicture}
			\end{center}
		\end{enumerate}

		\item \begin{enumerate}
			\item Sur l'intervalle $\left[ \dfrac43 ; +\infty \right[$, la fonction $f$ est strictement positive (elle est minorée par 2), donc l'équation $f(x)=0$ n'admet pas de solution.

			Sur l'intervalle $\left] -\infty  ;\dfrac43 \right]$, la fonction $f$ est dérivable donc continue. Elle est strictement croissante sur cet intervalle.

			De plus, l'image de l'intervalle $\left] -\infty  ;\dfrac43 \right]$ est l'intervalle $\left] -\infty  ;3\e^{\frac43}+2 \right]$ avec $3\e^{\frac43}+2 >0$.

			Donc d'après le théorème de la bijection, l'équation $f(x)=0$ admet une unique solution sur $\left] -\infty  ;\dfrac43 \right]$.

			En conclusion, l'équation $f(x)=0$ admet une unique solution $\alpha$ sur $\R$.

			\item On sait que $\alpha \in \left] -\infty  ;\dfrac43 \right]$

			La calculatrice permet d'obtenir :

			$f \left(-1 \right) \approx \np{-8,87}$ et $f \left(0 \right)=1$ donc $-1 < \alpha < 0 $

			$f \left( \np{-0,3} \right) \approx \np{-0,56}$ et $f \left( \np{-0,2} \right) \approx \np{-0,046}$ donc $- \np{0,3} < \alpha < \np{-0,2} $

			$f \left( \np{-0,21} \right) \approx \np{-0,01}$ et $f \left( \np{-0,2} \right) \approx \np{-0,046}$ donc $- \np{0,21} < \alpha < \np{-0,2} $

			$f \left( \np{-0,209} \right) \approx \np{-0,00}5$ et $f \left( \np{-0,208} \right) \approx \np{-0,00}05$ donc $- \np{0,209} < \alpha < \np{-0,208} $

			On en déduit que $\alpha \approx \np{-0,21}$. (Remarque : on peut se contenter de donner la valeur approchée renvoyée par le mode "résolution d'équation" de la calculatrice).

			\item Par lecture du tableau de variation, on obtient le tableau de signes suivant
			\begin{center}
				\begin{tikzpicture}
					\tkzTabInit[lgt=3,espcl=4]{$x$ / 0.8 , signe de $f$ / 1}{$-\infty$, $\alpha$, $+\infty$}
					\tkzTabLine{,-,z,+, }
				\end{tikzpicture}
			\end{center}

			\item Sur l'intervalle $\big[0  ; 2 \big]$, la fonction $f$ est positive donc l'aire du domaine limité par la courbe $\mathcal{C}_f$, l'axe de abscisses et les droites d'équation $x = 0$ et $x = 2$ est $\mathcal{A}=\displaystyle \int_0^2 f(x) dx$.

			$\begin{aligned}[t] \mathcal{A}&=\displaystyle \int_0^2 f(x) dx \\
			&= \bigg[ F(x) \bigg]_0^2 \\
			&= F(2) - F(0) \\
			&= -8\e^2 +4 +2 \\
			&=6-8\e^2
			\end{aligned}$

			L'aire du domaine limité par la courbe $\mathcal{C}_f$, l'axe de abscisses et les droites d'équation $x = 0$ et $x = 2$ mesure $6-8\e^2$ unités d'aire.
		\end{enumerate}

		\item On remarque que l'inéquation est :\quad $f(x) \leqslant 4x + 1$.

		Graphiquement, cela correspond au fait que la courbe $\mathcal{C}_f$ est en-dessous de sa tangente au point d'abscisse 0, et ce pour tout réel $x$ : cela fait penser à étudier la convexité de la fonction $f$.

		On va donc déterminer la dérivée seconde de $f$ :

		On a, pour tout $x$ réel :\quad  $f''(x)= -3\e^{-x} + (-3x + 4)\times \big(-\e^{-x}\big) = \e^{-x} \big(-3 - (-3x+4)\big)-\e^{-x} (3x-7)$.

		La dérivée seconde $f''$ est strictement négative quand $3x-7$ est strictement négatif soit sur $\left] -\infty  ; \dfrac73 \right[$.

		On en déduit que sur $\left] -\infty  ; \dfrac73 \right[$, $f$ est concave donc sa courbe $\mathcal{C}_f$ est en-dessous de toutes ses tangentes. En particulier, $\mathcal{C}_f$ est en-dessous de sa tangente au point d'abscisse 0.

		On a donc, pour tout $x \in \left] -\infty  ; \dfrac73 \right[$, \quad $(3x - 1)\e^{-x} + 2 \leqslant 4x + 1$.

		Sur $\bigg[ \dfrac73  ; +\infty \bigg[$, la fonction \quad $x \longmapsto 4x+1$ \quad est croissante donc \quad $4x+1 \geqslant 4 \times \dfrac73 +1$\quad soit \quad $4x+1 \geqslant \dfrac{29}{3}$.

		Sur $\bigg[ \dfrac73  ; +\infty \bigg[$, la fonction $f$ est décroissante donc \quad $f(x) \leqslant f\left( \dfrac73 \right)$ \quad or $f\left( \dfrac73 \right) \leqslant \dfrac{29}{3}$, \quad d'après la concavité de $f$ sur $\left] -\infty  ; \dfrac73 \right[$.

		On en conclut que pour tout $x \in \bigg[ \dfrac73  ; +\infty \bigg[$, \quad $f(x) \leqslant 4x+1$.
	\end{enumerate}

\needspace{10\baselineskip}

\end{document}
